\section{Metodología}

Para analizar el comportamiento de un contactor electromagnético, tanto en su estado de reposo como energizado se siguió el siguiente procedimiento:

\begin{enumerate}
    \item Desarmar el contactor para identificar sus componentes internos.
    \begin{itemize}
        \item Con mucho cuidado se desarmó el contactor y seguidamente se procedió a identificar y separar sus partes como la bobina, núcleo fijo, armadura o martillo, contactos principales, contactos auxiliares, espiras sombra y el resorte.
        \item Se observó el entrehierro permanente entre el núcleo fijo y la armadura, el cual es crucial para el funcionamiento del contactor.
        \item Se respondieron las preguntas planteadas para esta sección de la práctica.
    \end{itemize}
    
    \item Rearmar el contactor:\\ 
    Se volvió a armar el contactor asegurando que todos los contactos y resortes estaban correctamente colocados.
    
    \item Luego se procedió al montaje del circuito eléctrico.\\
    El cual consistió en conectar al contactor a una fuente de alimentación, luego a un auto transformador (variac) y luego a un transformador elevador para controlar el voltaje de entrada, y los instrumentos de medición necesarios para registrar las variables eléctricas durante las pruebas.
    
    \item Para medir la corriente de llamada se procedió a:
    \begin{itemize}
        \item Cerrar el interruptor de entrada para energizar el circuito.
        \item Luego aumentar gradualmente el voltaje de entrada hasta que el contactor cerrara.
        \item Tomar nota de la corriente antes de que el contactor cerrara, esta es conocida como la corriente de llamada.
    \end{itemize}

    \item Para medir la corriente de mantenimiento se procedió a:
    \begin{itemize}
        \item Colocar un transformador de corriente ya que la corriente de mantenimiento es muy pequeña para poder medirse con los instrumentos presentes en el laboratorio.
        \item Subir nuevamente el voltaje de entrada hasta que el contactor cerró. 
        \item Mantener el contactor energizado después del cierre.
        \item Tomar nota de la corriente mientras el contactor se mantiene cerrado, esta es conocida como la corriente de mantenimiento.
    \end{itemize}

    \item Para medir la tensión de mantenimiento y mínima operación se procedió a:
    \begin{itemize}
        \item Energizar el circuito y subir gradualmente el voltaje hasta que el contactor cerró.
        \item Se tomó nota de este voltaje que fue el voltaje mínimo de operaciones.
        \item Luego se disminuyó gradualmente el voltaje de entrada hasta que el contactor se abrió nuevamente.
        \item Se tomó nota del voltaje justo antes de que el contactor se abriera, esta es conocida como la tensión de mantenimiento.
    \end{itemize}

    \item Luego se tomó un segundo contactor para el cual se repitieron los pasos anteriores para hallar las corrientes de mantenimiento y llamada así como también las tensiones de mantenimiento y mínima de operaciones.
    \item Por último se observo que ocurre cuando al contactor se le colocó una hoja de papel entre el contacto metálico, en el cual se pudo escuchar un ruido muy fuerte y se procedió a tomar nota y registrar lo observado.
\end{enumerate}
